\begin{abstract}
比特币是目前市值最高的加密货币，开发工作完全依赖全球开源开发者的自愿贡献。
对激励开发者持续参与，始终是比特币社区及众多开源项目面临的核心挑战。
既有研究多关注实际货币激励对开发者行为的影响，
但捐赠信息披露本身所产生的激励效应，以及
社区内部情感状态对开发者活动的动态作用，至今缺乏系统考察。

本文以 GitHub 上的比特币开源项目为研究对象，围绕两个相互补充的问题展开：

其一，探究外部非盈利基金会的捐赠信息披露对开发者活动的影响。
Brink 运营中存在两次信息披露：在 X 平台公布收到的捐款信息（非定向披露），
以及在官网公布受资助开发者名单（定向披露）。
本文采用面板数据固定效应模型与工具变量方法，
剔除实际受资助开发者后识别纯信息效应。
结果表明，两种披露均对开发者活动产生显著正向影响，
但定向披露对非定向披露的正向效应形成负向调节——
这与期望失验理论所预测的失望效应相一致。

其二，探究比特币社区 Issue 讨论中的情感状态对开发者活动的影响，
并检验代码活动与讨论活动之间是否存在替代效应。
基于向量自回归模型的分析发现，滞后一期的负面情感对代码活动与讨论活动均有正向影响，
而滞后两期的代码活动则负向影响讨论活动，
由此揭示出“情感→代码/讨论→讨论”的传导路径。
上述结果分别借助情感即信息理论和目标手段替代理论予以解释。

本研究从信息视角重新审视开源开发者的激励机制，
为比特币及更广泛开源社区的捐赠信息披露策略设计提供实证参考，
同时对图书情报学中信息行为与用户激励研究作出理论贡献。
\end{abstract}

% 英文摘要
\begin{abstract*}
Bitcoin, the world's highest-valued cryptocurrency, relies entirely on voluntary contributions
from open-source Developmenters worldwide. Sustaining Developmenter motivation remains a central
challenge for the Bitcoin community and open-source projects at large. While prior research
has predominantly examined the direct effects of monetary incentives, the motivating impact
of \textit{donation information disclosure} itself—decoupled from actual financial transfers—
and the role of \textit{community sentiment} in shaping Developmenter behavior remain underexplored.

This thesis examines two complementary research questions using the Bitcoin GitHub repository
as the empirical context. First, it investigates how two distinct donation disclosure
strategies employed by the non-profit organization Bitcoin Brink affect Developmenter activity:
untargeted announcements of received donations posted on platform X, and targeted
grantee disclosures published on Brink's official website. Employing panel fixed-effects
models and instrumental variable estimation—while excluding actual grantees to isolate
pure information effects—the study finds that both disclosure types positively influence
Developmenter contributions. However, targeted disclosure negatively moderates the positive
effect of untargeted disclosure, consistent with the disappointment mechanism predicted
by expectancy disconfirmation theory.

Second, the thesis examines how community sentiment expressed in Issue discussions
dynamically influences Developmenter activities, and whether substitution effects exist between
code and discussion activities. Using Vector Autoregression (VAR) models, the analysis
reveals that lagged negative sentiment positively influences both code and discussion
activities, while increased code activity subsequently suppresses discussion activity.
This yields the influence pathway: sentiment → code/discussion → discussion, explained
through affect-as-information theory and goal–means substitution theory.

This research reframes open-source Developmenter motivation from an informational perspective,
offering actionable insights for donation disclosure strategy design in Bitcoin and
broader open-source communities, and contributing theoretically to information behavior
research in library and information science.
\end{abstract*}
